\chapter*{Use of Generative AI}
\addcontentsline{toc}{chapter}{Use of Generative AI}
\markboth{Use of Generative AI}{Use of Generative AI}

In keeping with current best practices for open educational resources,
this chapter discloses how generative artificial intelligence (AI) was
used in the preparation of this book.

%-------------------------------------------------------------------
\section*{Authorship}
%-------------------------------------------------------------------

The mathematical content, structure, examples, exercises, proofs, and
pedagogical decisions in this book are the work of Robert Hildebrand
and the contributors named in the Acknowledgements and Sources and
Attribution sections. AI was used as a writing-assistance and
production tool, not as an author. The author has reviewed, verified,
and edited all material that AI assisted with.

%-------------------------------------------------------------------
\section*{What AI Was Used For}
%-------------------------------------------------------------------

Generative AI assistants (primarily large language models) were used
during preparation of this book for the following kinds of tasks:

\begin{itemize}
  \item \textbf{Editing and copy-editing} of expository prose for
    clarity, grammar, and consistency of voice.
  \item \textbf{Drafting boilerplate text} such as license-compatibility
    summaries, attribution notices, accessibility metadata, and
    chapter-opening framing paragraphs, which the author then revised.
  \item \textbf{Generating draft alt-text descriptions} for figures,
    which the author reviewed and edited for mathematical accuracy.
  \item \textbf{Formatting assistance with \LaTeX{}}: generating
    \texttt{tikz} skeletons, fixing macro definitions, restructuring
    tables, debugging compilation errors.
  \item \textbf{Code-comment and example-code drafting} in Python and
    Julia. All numerical examples and code were independently verified
    by running the code and checking results.
  \item \textbf{License-audit assistance}: cross-checking source
    license compatibility, drafting license notices, identifying
    figures or text passages that needed re-attribution or rewriting.
  \item \textbf{Proofreading and consistency checks} across notation
    used in different chapters.
\end{itemize}

%-------------------------------------------------------------------
\section*{What AI Was Not Used For}
%-------------------------------------------------------------------

Generative AI was not used to:

\begin{itemize}
  \item Generate or fabricate mathematical claims, proofs, theorems,
    or numerical results that were not independently verified by the
    author.
  \item Produce final figures used in the published book without
    review. All TikZ figures were authored by hand or generated by AI
    and then reviewed and corrected.
  \item Generate exercise solutions without verification.
  \item Make pedagogical decisions about scope, sequencing, or
    emphasis. Those decisions are the author's.
\end{itemize}

%-------------------------------------------------------------------
\section*{Example Prompt Types}
%-------------------------------------------------------------------

To give readers and reviewers a concrete sense of how AI was engaged,
the following are representative prompt patterns that were used during
preparation:

\begin{itemize}
  \item ``Rewrite this paragraph to improve clarity without changing
    its mathematical content.''
  \item ``Generate a one-sentence alt-text description of the
    following TikZ figure that conveys both its visual layout and its
    mathematical meaning.''
  \item ``Verify that this CC~BY~3.0 source can be incorporated into a
    CC~BY-SA~4.0 work and explain the compatibility mechanism.''
  \item ``Draft a Python example using the \texttt{pyomo} package that
    solves the following linear program; do not invent solver options
    you are not certain exist.''
  \item ``Identify any LaTeX macros referenced but not defined in this
    file and propose corrections.''
\end{itemize}

%-------------------------------------------------------------------
\section*{Verification and Accountability}
%-------------------------------------------------------------------

All AI-assisted contributions were reviewed by the author. Numerical
examples were verified by running code and, where appropriate, by
independent computation in tools such as \texttt{scipy.optimize} and
\texttt{Gurobi}. The mathematical correctness of every statement,
proof, and example in this book is the author's responsibility.

Readers who suspect an error -- whether it was AI-introduced or not --
are warmly encouraged to report it via the project's GitHub issue
tracker or by emailing \texttt{open.optimization@gmail.com}. Errata
will be acknowledged in the book's revision history.

%-------------------------------------------------------------------
\section*{Disclosure Standard}
%-------------------------------------------------------------------

This disclosure follows the principles set out by Virginia Tech
Publishing's open educational resources program and is consistent with
emerging norms across the OER community: AI assistance is acceptable
in production of open educational materials when (i) the human
author remains responsible for content and accuracy, (ii) AI is not
listed as an author, and (iii) the use of AI is disclosed
transparently in the published work.

\cleardoublepage
